\documentclass[11pt]{article}
\usepackage{amsmath,amssymb,amsthm,booktabs,longtable,microtype}
\usepackage[a4paper,margin=26mm]{geometry}
\usepackage[colorlinks=true,linkcolor=blue,citecolor=blue,urlcolor=blue]{hyperref}

\newtheorem{theorem}{Theorem}
\newtheorem{proposition}[theorem]{Proposition}
\newtheorem{lemma}[theorem]{Lemma}
\newtheorem{corollary}[theorem]{Corollary}
\newtheorem{remark}[theorem]{Remark}

\newcommand{\Q}{\mathbf Q}
\newcommand{\F}{\mathbf F}
\newcommand{\Sel}{\operatorname{Sel}}

\title{A Two-Isogeny Descent Attached to Campbell's Eight-Term Family}
\author{Draft}
\date{Revised 4 September 2026}

\begin{document}
\maketitle

\begin{abstract}
Starting from Campbell's cubic Weierstrass family with eight rational points
whose $x$-coordinates are $0,\ldots,7$, we reconstruct the quartic condition
for the ninth candidate and record every rational specialization boundary.
We identify its Jacobian and rational $2$-torsion, then perform an exact
finite local calculation for the dual $2$-isogeny descents.  The resulting
Selmer groups have dimensions $3$ and $2$, respectively, which gives
Mordell--Weil rank at most $3$.  An explicit integral change gives a global
minimal semistable model of conductor $301245307115205810$.  We also compute the full $\Q\times K$
representative of the binary-quartic cubic invariant and its rational
$2$-torsion projection.  A versioned supplement records the $512$ local
cells, common-parameter local witnesses, exact generators, certificates, and
regression tests.  Uniform valuation arguments show that the $E'$-side
classes surviving both $\Q_2$ and $\Q_3$ are exactly the four Selmer
classes; together with the analogous $E$-side two-place gate, this replaces
the decisive portions of the local table by proofs.  The same-parameter genus-$5$ fibre product is everywhere
locally soluble.  Our calculation is deliberately narrower than the original
extension problem: it neither produces nor excludes a ninth rational point
on Campbell's family, claims no full $2$-Selmer group or Mordell--Weil rank
equality, and does not evaluate a Cassels--Tate pairing.  The results
are therefore a reproducible finite descent computation and a rigorously
delimited input to further arithmetic study of the ninth-value curve.
\end{abstract}

\noindent\textbf{Keywords.} elliptic curves; $2$-isogeny descent; Selmer
groups; arithmetic progressions; local solubility.

\medskip
\noindent\textbf{Mathematics Subject Classification (2020).} Primary 11G05;
Secondary 11G30, 11Y50.

\section{The Campbell quartic}
Campbell \cite[Corollary~2.4 and Theorem~2.5]{Campbell2003} uses the
parameter change
\[
 t=\frac{6m^2-126m-285360}{m^2-72256}.
\]
Its denominator never vanishes for $m\in\Q$, since
$268^2<72256<269^2$.  In the normalization used here put
\[
 g_m(x)=g_3x^3+g_2x^2+g_1x+g_0,
\]
where
\begin{align*}
g_3={}&-18816m^4+677376m^3+1922543616m^2-48944480256m
-40678301368320,\\
g_2={}&236896m^4-9821952m^3-22598349824m^2+508953231360m
+520252184657920,\\
g_1={}&-958800m^4+40985280m^3+89932669440m^2-1957723729920m
-2113363439616000,\\
g_0={}&1292769m^4-57304800m^3-118795148928m^2+2647001548800m
+2758336954896384.
\end{align*}
Exact expansion gives $g_m(j)=q_j(m)^2$ for $0\le j\le6$; the square roots
are listed in Table~\ref{tab:campbell-roots}.
\begin{table}[ht]
\centering\small
\caption{The first seven prescribed square roots in Campbell's family}
\label{tab:campbell-roots}
\begin{tabular}{c@{\quad}l}
\toprule
$j$&$q_j(m)$\\ \midrule
0&$3(379m^2-8400m-17506624)$\\
1&$743m^2-17136m-33534272$\\
2&$415m^2-11088m-16946752$\\
3&$3(67m^2-3696m-494656)$\\
4&$209m^2-17136m+5050432$\\
5&$263m^2-25200m+10631872$\\
6&$63(m^2-2352m+72256)$\\ \bottomrule
\end{tabular}
\end{table}
The eighth condition is
\[
\begin{aligned}
Y^2=D(m)={}&-264815m^4-19343520m^3+62846856064m^2\\
&-2906312951808m-495507443511296,
\end{aligned}
\]
because $g_m(7)=D(m)$.  The candidate immediately following the eight
indices $0,\ldots,7$ has index $8$, and direct expansion gives $g_m(8)=H(m)$,
where
\[
 \begin{aligned}
 H(m)={}&-850079m^4-11210976m^3+138714149248m^2\\
        &-5501355374592m-1679721044504576
 \end{aligned}
\]
and put $C_H:z^2=H(m)$.  Thus Campbell's construction supplies eight
prescribed squares precisely when $(m,Y)$ lies on $Y^2=D(m)$, and
$z^2=H(m)$ tests the ninth candidate at the same $m$.

There are no omitted rational specialization boundaries.  The leading
coefficient factors as
\[
g_3=-18816(m^2-72256)(m^2-36m-29920),
\]
and its two quadratic factors have no rational roots.  After division by
$-65028096$, the discriminant of $g_m$ in $x$ is a primitive degree-$16$
polynomial irreducible modulo $53$, hence it has no rational zero.  The point
$m=\infty$ is not rational on $Y^2=D(m)$, since its leading equation would be
$Y^2=-264815M^4$.  A zero of $D$ or $H$ merely gives $Y=0$ or $z=0$ at one
of the distinct $x$-coordinates; it does not make $g_m$ singular.  Finally,
$D$ and $H$ are separable and coprime, so their smooth fibre product has
genus $5$.

\begin{lemma}[same-parameter local solubility]\label{lem:same-m}
The smooth curve
\[
Y^2=D(m),\qquad z^2=H(m)
\]
has points over $\mathbb R$ and every $\Q_p$.
\end{lemma}
\begin{proof}
The supplement stores a common $m$ and both square roots at $\mathbb R$,
$\Q_2$, every odd prime below $101$, and every odd prime dividing the two
quartic discriminants or their resultant.  The branch-bad set is
\[
\{2,3,5,7,17,19,31,59,8599,71699,898543,23037169,
339106321,1153266911\}.
\]
At every remaining prime $p\ge101$ the fibre product has smooth genus-$5$
reduction, and the Weil bound
$\#C(\F_p)\ge p+1-10\sqrt p>0$ supplies a smooth point, which lifts by
Hensel's lemma.  The exact witness table is the object
\texttt{same\_m\_local\_certificates} in the manifest-bound certificate.
\end{proof}
In particular $C_H$ is everywhere locally soluble.  Nothing here asserts
that the fibre product or $C_H$ has a rational point.

For a binary quartic
\[
 g(X,Z)=a_0X^4+a_1X^3Z+a_2X^2Z^2+a_3XZ^3+a_4Z^4
\]
we use
\[
 I=12a_0a_4-3a_1a_3+a_2^2,
 \quad
 J=72a_0a_2a_4-27a_0a_3^2-27a_1^2a_4
      +9a_1a_2a_3-2a_2^3.
\]
The resulting Jacobian is $y^2=x^3-27Ix-27J$; see
\cite[Lemma~2.1]{Fisher2022} for the binary-quartic description of the
full $2$-Selmer group.

\section{The cubic algebra and the Campbell class}
Let $L=\Q[\varphi]/(\varphi^3-3I\varphi+J)$.  For the above normalization,
the cubic invariant of $g$ is
\begin{equation}\label{eq:cubic-invariant}
 z(g)=\frac{4a_0\varphi+3a_1^2-8a_0a_2}{3}
       \in L^\times/L^{\times2};
\end{equation}
this is the invariant used in \cite[Section~2]{Fisher2022}.  In the present
case
\[
 L\simeq\Q\times K,\qquad
K=\Q(w),\quad w^2=D,\quad
 D=1434501462453361=59\cdot71699\cdot339106321.
\]
The quadratic resolvent factor is
\[
 T^2+269378023424T-36009487121810563530752,
\]
whose discriminant is $12288^2D$.  Thus its roots are
$-134689011712\pm6144w$.  Although $\mathcal O_K=\mathbb
Z[(1+w)/2]$, all displayed coordinates below use the $\Q$-basis $(1,w)$.

\begin{proposition}\label{prop:z}
The two components of the cubic invariant of $H$ are
\[
 z_{\Q}=9250179026780160=35\cdot16257024^2
\]
and
\[
 z_K=467235380575281152-6963847168w
     =64^2(114071137835762-1700158w).
\]
Moreover
\[
 N_{K/\Q}(114071137835762-1700158w)
   =35\cdot15915620907648^2,
\]
and the full norm is
\[
 z_{\Q}N_{K/\Q}(z_K)
  =37093056870271943410974720^2.
\]
\end{proposition}
\begin{proof}
Insert the five coefficients of $H$ into \eqref{eq:cubic-invariant}, then
evaluate in the rational and quadratic factors of $L$.  Squaring the three
displayed integers gives the claimed identities.  These are integer
identities, independently recomputed by the accompanying regression test.
\end{proof}

The rational factor in $L=\Q\times K$ is the cohomological factor induced by
the rational $2$-torsion point.  On scaling
\[
 x_{\rm big}=64^2x,\qquad y_{\rm big}=64^3y
\]
and translating $X=x+197298357$, the Jacobian becomes
\begin{equation}\label{eq:E}
 E:y^2=X^3-591895071X^2+58536289153843200X.
\end{equation}
Since $-3\varphi_0=64^2(-197298357)$ for the rational resolvent root
$\varphi_0=269378023424$, the rational projection of $[H]\in
H^1(\Q,E[2])$ is the $X$-Kummer squareclass $[35]$.  This assertion concerns
one projection only: $C_H$ is not thereby identified with the even quartic
$C_{35}$ below.

\subsection{A global minimal model and conductor}

The model \eqref{eq:E} is convenient for the rational $2$-isogeny but is
not globally minimal.  The admissible change
\begin{equation}\label{eq:minimal-change}
 X=36x,\qquad y=216y_0+108x
\end{equation}
gives
\begin{equation}\label{eq:Eminimal}
 E_{\min}:\quad y_0^2+xy_0
 =x^3-16441530x^2+45166889779200x.
\end{equation}
Equivalently, the admissible parameters are $(u,r,s,t)=(6,0,3,0)$, and the
$a$-invariants of \eqref{eq:Eminimal} are
$(1,-16441530,0,45166889779200,0)$.

\begin{proposition}\label{prop:minimal-model}
Equation \eqref{eq:Eminimal} is a global minimal model.  Its invariants are
\[
 c_4=2157171698920561,\qquad
 c_6=70577985500751764077559,
\]
and
\[
 \begin{split}
 \Delta_{\min}
 &=2926451742397178075653974744686961623040000\\
 &=2^{28}3^{16}5^4 7^{10}\cdot59\cdot71699\cdot339106321.
 \end{split}
\]
The curve is semistable, with multiplicative reduction of type
$I_{v_p(\Delta_{\min})}$ at each prime in
\[
\{2,3,5,7,59,71699,339106321\},
\]
and good reduction elsewhere.  The multiplicative reduction at $2$ and at
$3$ is split.  In particular,
\[
 N_E=2\cdot3\cdot5\cdot7\cdot59\cdot71699\cdot339106321
     =301245307115205810.
\]
Finally,
\[
 j(E)=
 \frac{10038160648206649953061393462836377818780518481}
 {2926451742397178075653974744686961623040000}.
\]
\end{proposition}
\begin{proof}
Substitution in the general admissible-change formulas gives
\eqref{eq:Eminimal}.  Direct calculation gives the displayed $c_4,c_6$
and discriminant and verifies $c_4^3-c_6^2=1728\Delta_{\min}$; the original
invariants are respectively $6^4,6^6,6^{12}$ times these invariants.
Moreover
\[
 \gcd(c_4,\Delta_{\min})=1.
\]
At every prime dividing $\Delta_{\min}$ we therefore have
$v_p(c_4)=0$.  An integral Weierstrass equation that were nonminimal at
$p$ would have $p^4\mid c_4$, so this proves $p$-minimality.  This argument
is valid without change at $p=2$ and $p=3$.

For an explicit check at these two residue characteristics, move the right
side of \eqref{eq:Eminimal} to the left.  Both $16441530$ and
$45166889779200$ vanish modulo $2$ and modulo $3$, so the reduced affine
equation is
\[
 F(x,y)=y^2+xy-x^3=0.
\]
Modulo $2$ its partial derivatives are $(F_x,F_y)=(y+x^2,x)$, and modulo
$3$ they are $(F_x,F_y)=(y,2y+x)$.  Thus in either case the unique affine
singular point is $(0,0)$; the point at infinity is smooth because the
$Z$-derivative of the homogenized equation is nonzero there.  At $(0,0)$
the quadratic tangent cone is
\[
 y^2+xy=y(y+x),
\]
whose two linear factors are distinct over both $\mathbf F_2$ and
$\mathbf F_3$.  Hence the multiplicative reduction at $2$ and $3$ is
split.  The criterion $v_p(\Delta)>0$, $v_p(c_4)=0$ gives multiplicative
reduction at every other bad prime as well, with conductor exponent one.
At every other prime the discriminant is a unit.  The conductor and Kodaira
symbols follow, and $j=c_4^3/\Delta_{\min}$.  See
\cite[Chapter~VII, Section~5]{Silverman2009} for the reduction criteria.
\end{proof}

This calculation supplies exact isomorphism-class search keys.  It is not a
rank computation, and the absence of a database row at this large composite
conductor is not evidence for novelty.

\section{The rational two-isogeny}
Write $A=-591895071$ and $B=58536289153843200$.  The quotient by
$(0,0)$ gives
\[
 \phi:E\longrightarrow E',\qquad
 E':y^2=X^3-2AX^2+(A^2-4B)X,
\]
so explicitly
\begin{equation}\label{eq:Eprime}
 E':y^2=X^3+1183790142X^2+116194618458722241X.
\end{equation}
We denote the dual isogeny by $\widehat\phi:E'\to E$.  The standard
isogeny descent and its covering quartics are described, for example, in
\cite[Chapter~X, Section~4]{Silverman2009} and
\cite[Chapter~14]{Cassels1991}; higher descent with rational $2$-torsion is
treated directly in \cite{Fisher2017}.

\begin{lemma}[support]\label{lem:support}
For $E_{a,b}:y^2=x^3+ax^2+bx$, every class in the $x$-Kummer isogeny
Selmer group has a squarefree representative supported on $-1$ and the
prime divisors of $b$.  The class $d$ is represented by
\begin{equation}\label{eq:Cd}
 C_d:\quad N^2=dU^4+aU^2V^2+(b/d)V^4.
\end{equation}
\end{lemma}
	\begin{proof}
Let $p\nmid b$ divide a squarefree representative $d$, and after weighted
projective scaling take $U,V\in\mathbb Z_p$ primitive.  In the original
equation \eqref{eq:Cd}, if $V$ were a unit, then $(b/d)V^4$ would be the
unique term of valuation $-1$ on the right, which cannot be the valuation
of the square $N^2$.  Hence $p\mid V$, so primitivity makes $U$ a unit.
The right side of \eqref{eq:Cd} then has valuation exactly $1$, again
impossible for a square.  Notice that this argument does not assume in
advance that $N$ is $p$-integral.  Thus $v_p(d)$ is even outside $b$, and
squareclass reduction gives the stated support.
	\end{proof}

\begin{proposition}[two-place $E$-side gate]\label{prop:two-place-gate}
For the curve $E$ in \eqref{eq:E}, let $d$ run through the $32$ signed
squarefree classes supported on $\{2,3,5,7\}$.  Then
\[
 C_d(\mathbf Q_{59})\ne\varnothing
 \quad\Longleftrightarrow\quad
 C_d(\mathbf Q_{71699})\ne\varnothing.
\]
These equivalent conditions hold exactly for
\[
 d\in\left\{
 \begin{gathered}
 -210,-70,-42,-30,-14,-10,-6,-2,\\
 1,3,5,7,15,21,35,105
 \end{gathered}
 \right\}.
\]
Moreover $C_d(\mathbf R)\ne\varnothing$ exactly when $d>0$.
Consequently the classes that survive the real place together with either
one of these two finite places are exactly
\[
 \{1,3,5,7,15,21,35,105\}.
\]
\end{proposition}
\begin{proof}
Write $a=-591895071$, $b=58536289153843200$, and
\[
 \delta=a^2-4b=3^4\cdot59\cdot71699\cdot339106321.
\]
For $p=59$ or $71699$ one has $v_p(\delta)=1$, while every listed $d$
and $2b$ are $p$-adic units.  The covering quartic satisfies
\begin{equation}\label{eq:double-root-identity}
 4dF_d=(2dU^2+aV^2)^2-\delta V^4,
 \qquad F_d=dU^4+aU^2V^2+(b/d)V^4.
\end{equation}
If $(d/p)=1$, the choice $(U:V)=(1:0)$ gives a point because $d$ is a
square in $\mathbf Q_p$.  Conversely suppose $(d/p)=-1$ and normalize
$U,V\in\mathbf Z_p$ to be primitive.  If $p\mid V$, then
$F_d\equiv dU^4\pmod p$ is a nonsquare unit.  If $V$ is a unit and
$2dU^2+aV^2$ is a unit, \eqref{eq:double-root-identity} shows that
$F_d$ again has squareclass $d$ modulo $p$.  In the remaining case,
$p\mid2dU^2+aV^2$; the right side of
\eqref{eq:double-root-identity} then has valuation exactly one, so
$v_p(F_d)=1$.  Each case contradicts $F_d=N^2$.

Quadratic reciprocity gives, at both primes,
\[
 \left(\frac{-1}{p}\right)=
 \left(\frac{2}{p}\right)=-1,\qquad
 \left(\frac{3}{p}\right)=
 \left(\frac{5}{p}\right)=
 \left(\frac{7}{p}\right)=1.
\]
This yields the displayed $16$-class local list.  Finally, if $d>0$ then
$(U:V)=(1:0)$ gives a real point.  If $d<0$, all three coefficients
$d,a,b/d$ of $F_d$ are negative, so $F_d<0$ for every nonzero real pair
$(U,V)$.  Intersecting the two classifications gives the eight positive
odd divisors of $105$.
\end{proof}

\begin{proposition}[two-place $E'$-side gate]\label{prop:eprime-two-place-gate}
Put
\[
 A'=1183790142,\qquad B'=116194618458722241=3^4D,
 \qquad q=339106321.
\]
For the $32$ signed squarefree classes $d$ supported on
$\{3,59,71699,q\}$, write
\[
 C'_d:\quad N^2=dU^4+A'U^2V^2+(B'/d)V^4.
\]
Then
\[
 C'_d(\mathbf Q_2)\ne\varnothing
 \quad\Longleftrightarrow\quad d\equiv1\pmod8,
\]
and its eight solutions among the supported classes are
\[
 \{1,q\}\{1,3\cdot59,3\cdot71699,59\cdot71699\}.
\]
Moreover
\[
 C'_d(\mathbf Q_3)\ne\varnothing
 \quad\Longleftrightarrow\quad v_3(d)=0\ \hbox{ and }\ d\equiv1\pmod3,
\]
and the corresponding eight classes are
\[
 \{1,q\}\{1,59\cdot71699,-59,-71699\}.
\]
Consequently simultaneous solubility at $2$ and $3$ leaves exactly
\[
 \{1,4230241,339106321,D\}.
\]
No single finite place in the stored eight-place matrix leaves only these
four classes; thus two is the minimum number of finite places that recovers
this four-class set from that matrix.  The only other minimal pair is
$\{2,5\}$.
\end{proposition}
\begin{proof}
Write $A'=2c$, so $c=591895071\equiv7\pmod8$.  Direct integer
factorization gives
\[
 B'=3^4\cdot59\cdot71699\cdot q\equiv1\pmod8,
 \qquad (A')^2-4B'=2^{22}k,
\]
where $k=223298222175$ is odd.  If $F'_d$ denotes the quartic on the
right of the equation for $C'_d$, coefficient comparison gives
\[
 4dF'_d=(2dU^2+A'V^2)^2-((A')^2-4B')V^4.
\]
Thus every point satisfies
\begin{equation}\label{eq:eprime-two-identity}
 dN^2=T^2-2^{20}kV^4,\qquad T=dU^2+cV^2.
\end{equation}

Normalize a putative $2$-adic point so that $U,V\in\mathbf Z_2$ are
primitive.  Suppose $d\not\equiv1\pmod8$.  When $U,V$ have opposite
parity, $T$ is odd.  When both are odd, their squares are $1$ modulo $8$,
and $T\equiv d+7\pmod8$; hence
\[
 v_2(T)=1,2,1\quad\hbox{as}\quad d\equiv3,5,7\pmod8.
\]
In every case $v_2(T)\le2$, and therefore
\[
 v_2\left(2^{20}kV^4/T^2\right)\ge16.
\]
Every element of $1+8\mathbf Z_2$ is a square (equivalently, apply the
strengthened Hensel criterion at the approximate square root $1$).
Equation~\eqref{eq:eprime-two-identity} consequently makes $dN^2$ a
nonzero square in $\mathbf Q_2$, forcing $d$ itself to be a square, a
contradiction.  Conversely, if $d\equiv1\pmod8$, then $d$ is a square in
$\mathbf Q_2$: for $f(X)=X^2-d$, one has
$v_2(f(1))\ge3>2v_2(f'(1))=2$.  The point
$(U:V:N)=(1:0:\sqrt d)$ proves sufficiency.

For the $3$-adic assertion, note that $v_3(A')=2$,
$A'/9\equiv1\pmod3$, and $B'/3^4=D\equiv1\pmod3$.  Again normalize
$U,V\in\mathbf Z_3$ to be primitive.  If $v_3(d)=1$, then the three
terms of $F'_d$ show that $v_3(F'_d)=1$ when $U$ is a unit.  If $3\mid U$,
then $V$ is a unit and the last term is the unique term of least valuation,
giving $v_3(F'_d)=3$.  Neither valuation can occur for a square.

Now let $v_3(d)=0$.  If $U$ is a unit, reduction modulo $3$ gives
$F'_d\equiv dU^4$, so solubility requires $d\equiv1\pmod3$.  If
$U=3^ru$ with $r\ge1$ and $u,V$ units, divide by $3^4$.  For $r\ge2$ the
result is congruent to $(D/d)V^4\equiv d^{-1}\pmod3$.  For $r=1$ it is
congruent to
\[
 du^4+(A'/9)u^2V^2+(D/d)V^4\equiv d+1+d^{-1}\pmod3.
\]
When $d\equiv-1\pmod3$, both displayed units are $-1$, so no square is
possible.  Conversely $d\equiv1\pmod3$ is a square in $\mathbf Q_3$ by
ordinary Hensel lifting of the root $1$ of $X^2-d$ modulo $3$, and the
same point $(1:0:\sqrt d)$ works.

Finally $59\equiv71699\equiv3\pmod8$, $q\equiv1\pmod8$, while
$59\equiv71699\equiv-1\pmod3$ and $q\equiv1\pmod3$.  These residues give
the two factored eight-class lists and their four-class intersection.
The last minimality statement is an exact comparison with all seven finite
place columns of the $512$-cell certificate: the smallest one-place survivor
count is eight, and precisely the pairs $\{2,3\}$ and $\{2,5\}$ have the
displayed four-class intersection.  This comparison is independently
recomputed by the standard-library regression script; the local
classification itself is proved above without using the table.
\end{proof}

\begin{lemma}[good primes]\label{lem:good}
Fix a supported squarefree $d$.  If a finite prime $p$ does not divide
$2b(a^2-4b)$, then the smooth projective model of $C_d$ has a $\Q_p$-point.
\end{lemma}
\begin{proof}
The binary-quartic discriminant of the right side of \eqref{eq:Cd} is
\[
 16b(a^2-4b)^2.
\]
Thus $C_d$ has smooth geometrically connected genus-one reduction at $p$.
In this application such a $p$ is at least $5$, so the Hasse bound gives
\(
 \#C_d(\F_p)\ge p+1-2\sqrt p>0.
\)
Every point of the smooth reduction lifts by Hensel's lemma.
\end{proof}

\section{Exact local classification}
The only places required on either side are
\[
 S=\{\infty,2,3,5,7,59,71699,339106321\}.
\]
Lemma~\ref{lem:support} gives $32$ signed candidates on each side, and
Lemma~\ref{lem:good} proves that no unlisted place can remove one.  The
certificate contains all $64\times8=512$ local cells.  Propositions
\ref{prop:two-place-gate} and \ref{prop:eprime-two-place-gate} now give
uniform proofs for the decisive $E$- and $E'$-side reductions.  A YES cell contains
an exact squareclass witness; a NO cell is either an exhaustive
weighted-projective computation modulo the displayed prime power or the
valuation-normalized double-root obstruction.  All $56$ cells that had
previously been unresolved are now split into $24$ YES cells at $3$ and
$32$ NO cells at $59$ or $71699$.

\begin{table}[ht]
\centering
\caption{Exact local classification of the two isogeny-descent candidate sets}
\label{tab:local-classification}
\begin{tabular}{lrrrrl}
\toprule
side & candidates & first real NO & first finite NO & survivors & survivor set\\
\midrule
$E$ & 32 & 16 & $8$ at $59$ & 8 & $1,3,5,7,15,21,35,105$\\
$E'$& 32 & 0 & $24$ at $2$, $4$ at $3$ & 4 & $1,4230241,339106321,D$\\
\bottomrule
\end{tabular}
\end{table}
The complete matrix has $384$ YES and $128$ NO cells.  Its JSON SHA-256 is
\begin{quote}\small\ttfamily
fff2f35f398c2d14227d9b032d205e24d5332a39346487c76180cdf805ba32c9.
\end{quote}
The structured $\Q\times K$ and Selmer certificate has SHA-256
\begin{quote}\small\ttfamily
7fbd4d00435d2f3ef7bd74b0f87822b8e6fb3b893e2023ce9b5741c411538e18.
\end{quote}
This is the clean Round-04 certificate: it contains only the exact Selmer,
$\Q\times K$, scaling, and proved Mordell--Weil-image data.  The superseded
opposite-side pairing proposal is absent; its failure is recorded only in the
separate Round-05 negative audit.
Both files, the same-$m$ certificate, their generators and regression tests
are located by supplement release
\begin{center}
\path{paper-elliptic-campbell-supplement-v0.6.1}
\end{center}
with root manifest
\begin{center}
\path{PAPER_ELLIPTIC_SUPPLEMENT_MANIFEST.json}.
\end{center}
The manifest SHA-256 is
\begin{quote}\small\ttfamily
583ec02b13c96824df5563d0dc7cc0312b7fbb33e5b7860b4f692b1fbed228fd.
\end{quote}
The manifest records byte lengths, roles, evidence eligibility, Python and
SymPy versions, and exact reproduction commands.  The current working tree is
public at \url{https://github.com/lixiang90/vibemath}; its immutable archive
identifier remains to be supplied.

\begin{theorem}\label{thm:selmer}
With $E$-side classes represented by the $X$-Kummer map on $E$ and
$E'$-side classes by that on $E'$, one has
\[
 \Sel^{\widehat\phi}(E'/\Q)=\langle3,5,7\rangle,
 \qquad
 \Sel^{\phi}(E/\Q)=\langle4230241,339106321\rangle.
\]
Their $\F_2$-dimensions are respectively $3$ and $2$.
\end{theorem}
\begin{proof}
The two support lists are exhaustive by Lemma~\ref{lem:support}; all good
places are automatic by Lemma~\ref{lem:good}.  Proposition
\ref{prop:two-place-gate} removes every $E$-side class outside the first
displayed group, and Proposition~\ref{prop:eprime-two-place-gate} does the
same on the $E'$ side.  The exact witnesses in the certificate verify that
every displayed class is soluble at the remaining bad places, yielding
exactly the survivor sets in Table~\ref{tab:local-classification}.  Direct squareclass multiplication
verifies that the two sets are
the groups generated as displayed.
\end{proof}

The exact sequence underlying descent by a degree-$2$ isogeny gives
\[
 2^{\operatorname{rank}E(\Q)}=
 \frac{|E'(\Q)/\phi E(\Q)|\,|E(\Q)/\widehat\phi E'(\Q)|}
 {|E(\Q)[\phi]|\,|E'(\Q)[\widehat\phi]|};
\]
see \cite[pp.~368--369]{vanBeekFisher2018}.  Both denominator groups here
have order $2$, and the two Mordell--Weil quotients inject into the two
Selmer groups.
\begin{corollary}\label{cor:rank}
$\operatorname{rank}E(\Q)\le3$.
\end{corollary}

\section{Correction of the proposed pairing gate}
For clarity we record the relevant cohomological direction.  The standard
pairing that upgrades a descent by $\widehat\phi$ to a full $2$-descent is
an alternating pairing
\begin{equation}\label{eq:isog-pairing}
 \Sel^{\widehat\phi}(E'/\Q)\times
 \Sel^{\widehat\phi}(E'/\Q)\longrightarrow\Q/\mathbb Z,
\end{equation}
whose radical is the image of the appropriate full $2$-Selmer group; see
\cite[Equation~(2)]{vanBeekFisher2018}.  In particular, it is a pairing on
one isogeny Selmer group squared, not a pairing between the two groups in
Theorem~\ref{thm:selmer}.  The full Cassels--Tate pairing on binary quartics
likewise has domain $\Sel^{(2)}(E/\Q)\times\Sel^{(2)}(E/\Q)$
\cite[Theorem~3.1]{Fisher2022}.

\begin{proposition}[pairing correction]\label{prop:correction}
The expression formerly denoted
$\langle35,4230241\rangle_{CT}$ in this project is not defined by either
of the standard pairings just cited: $35$ and $4230241$ are coordinates in
opposite isogeny Selmer groups.  Consequently its proposed tangent-line
formula cannot be used to deduce $C_H(\Q)=\varnothing$.
\end{proposition}
\begin{proof}
The domain mismatch follows directly from \eqref{eq:isog-pairing} and
Theorem~\ref{thm:selmer}.  There is a second consistency check.  By
Lemma~\ref{lem:same-m}, $C_H$ is everywhere locally soluble, so it represents a class of
$\Sel^{(2)}(E/\Q)$, and its projection is $35$.  Hence $35$ already lies in
the image whose radical is tested by \eqref{eq:isog-pairing}; any correctly
defined same-side lifting pairing must annihilate it.
\end{proof}

The arithmetic that exposed the problem is still useful as a negative
certificate.  For the standard even quartic
\[
 N^2=35U^4-591895071U^2V^2+1672465404395520V^4
\]
the associated conic in $R=U^2,S=V^2$ has the rational point
\[
 (R_0,S_0,N_0)=(16257024,1,36058176).
\]
A primitive tangent there is
\[
 L_{35}=60677401R-697502396215296S-8012928N.
\]
Both identities follow by direct substitution.  If one nevertheless
evaluates the rejected expression $(L_{35}(P_v),4230241)_v$, then the two
Hensel branches give the exact residues in
Table~\ref{tab:rejected-hilbert}.

\begin{table}[ht]
\centering
\caption{Branch dependence of the rejected local Hilbert-symbol expression}
\label{tab:rejected-hilbert}
\begin{tabular}{rrrrrr}
\toprule
$p$ & modulus & $N$ & $L_{35}$ & $v_p(L_{35})$ & Hilbert symbol\\
\midrule
$59$ & $59^3$ & $56226$ & $127303$ & $0$ & $-1$\\
$59$ & $59^3$ & $149153$ & $145730$ & $1$ & $+1$\\
$71699$ & $71699^3$ & $152989401412805$ & $252464922339130$ & $0$ & $-1$\\
$71699$ & $71699^3$ & $215596989132294$ & $195462237955773$ & $1$ & $+1$\\
\bottomrule
\end{tabular}
\end{table}
Each listed $N$ squares to the quartic right side modulo the listed modulus.
Thus independent local branch choices yield both possible products.  This
does not compute a Cassels--Tate value; it proves that the formerly proposed
bare expression lacks the cochain or denominator data needed for
well-definedness.

\section{The exact remaining full-descent problem}
To test the Campbell class itself one must now do one of the following
equivalent full-descent computations.
\begin{enumerate}
\item Compute a basis of $\Sel^{(2)}(E/\Q)$ as everywhere locally soluble
binary quartics $g_i$ with the same $I,J$; identify $[H]$ by the full
$z(H)\in(\Q\times K)^\times/(\Q\times K)^{\times2}$; and compute
$\langle[H],[g_i]\rangle_{CT}$ by the full binary-quartic formula of
\cite[Theorem~3.1]{Fisher2022}.
\item Compute the locally soluble $4$-coverings lying above $C_H$.  An empty
output proves that its $2$-Selmer class does not lift and hence that
$C_H(\Q)=\varnothing$; a nonempty output does not by itself produce a
rational point.  This is the precise role of a four-descent
\cite[Section ``Four-Descent'']{MagmaHandbook}.
\end{enumerate}
The accompanying Magma input freezes both routes, but the manifest marks it
\path{UNEXECUTED_FROZEN_CANDIDATE_INPUT} and
\path{mathematical_evidence_eligible=false}.  No transcript or trusted
binary hash is present.  Therefore no full $2$-Selmer dimension,
Cassels--Tate value, Mordell--Weil rank equality, or rational-point conclusion
is claimed.

\section{Finite theorem and reproducibility boundary}
\begin{theorem}[finite Campbell--Jacobian calculation]
For the binary quartic $H$ obtained at index $8$ from Campbell's
Theorem~2.5, the two isogeny Selmer groups are those in
Theorem~\ref{thm:selmer}; hence $\operatorname{rank}E(\Q)\le3$.  Its
cubic-algebra invariant is the pair of Proposition~\ref{prop:z}, with
rational $2$-torsion projection $[35]$.  Its global minimal model,
discriminant, local reduction types, and conductor are those of
Proposition~\ref{prop:minimal-model}.  Moreover the same-parameter
genus-$5$ fibre product is everywhere locally soluble.
\end{theorem}
This is a finite, certificate-backed theorem.  Exact-coefficient searches did
not locate a published computation of these particular Selmer groups; that is
a not-found report, not proof of novelty or priority.  The theorem neither
asserts the existence or nonexistence of a rational ninth point nor computes
the full $2$-Selmer group or any Cassels--Tate value.  No independent second
elliptic-curve CAS was available; the exact standard-library local proof and
compatibility test are not represented as an independent external
reproduction.  In particular, the local theorem does not decide the
ninth-point problem.

\section*{Statements and Declarations}

\paragraph{Funding.}
\emph{Placeholder: the human authors must insert the verified funding
statement.}

\paragraph{Competing interests.}
\emph{Placeholder: the human authors must insert the verified financial and
non-financial interests statement.}

\paragraph{Author contributions.}
\emph{Placeholder: the human authors must insert the approved contribution
statement.}

\paragraph{Ethics approval, consent to participate, and consent for publication.}
Not applicable: this article contains only mathematical research and uses no
human participants, personal data, or animals.

\paragraph{Data and code availability.}
The exact generators, JSON certificates, and regression tests are identified
by supplement release
\path{paper-elliptic-campbell-supplement-v0.6.1}.  The current working version
is public at \url{https://github.com/lixiang90/vibemath}.  Before submission,
the authors must freeze those exact bytes and provide a verified persistent
DOI or immutable release URL: \texttt{<<PERSISTENT DOI/URL REQUIRED>>}.

\paragraph{Use of generative AI and AI-assisted technologies.}
AI-assisted tools were used during exploratory symbolic-algebra scripting,
test drafting, literature-query formulation, and prose preparation.  No AI
system is an author.  Before submission, the named human authors must
independently verify and accept every mathematical claim, reference, source
file, code artifact, and disclosure, and they remain accountable for the
final work.

\paragraph{Supplementary information.}
Online Resource 1: deterministic verification archive containing the exact
source, frozen reference PDF, generators, certificates, data dictionary, and regression tests
identified by release root
\path{paper-elliptic-campbell-local-release-v0.6.3}.  The unexecuted Magma
input is included solely as an ineligible candidate input and is not evidence
for any theorem.  The journal name, author names, affiliations, corresponding
author email, and public archive identifier remain human-supplied metadata.

{\small
\begin{thebibliography}{9}
\bibitem{Campbell2003}
G. Campbell,
``A note on arithmetic progressions on elliptic curves,''
\emph{Journal of Integer Sequences} \textbf{6} (2003), Article 03.1.3,
\url{https://cs.uwaterloo.ca/journals/JIS/VOL6/Campbell/campbell4.html}.

\bibitem{Cassels1991}
J. W. S. Cassels,
\emph{Lectures on Elliptic Curves}, LMS Student Texts 24,
Cambridge University Press, 1991.

\bibitem{Fisher2022}
T. Fisher,
``On binary quartics and the Cassels--Tate pairing,''
\emph{Research in Number Theory} \textbf{8} (2022), article 74,
\url{https://doi.org/10.1007/s40993-022-00376-z}.

\bibitem{Fisher2017}
T. Fisher,
``Higher descents on an elliptic curve with a rational $2$-torsion point,''
\emph{Mathematics of Computation} \textbf{86} (2017), 2493--2518,
\url{https://doi.org/10.1090/mcom/3163}.

\bibitem{Silverman2009}
J. H. Silverman,
\emph{The Arithmetic of Elliptic Curves}, second edition,
Graduate Texts in Mathematics 106, Springer, 2009.

\bibitem{vanBeekFisher2018}
M. van Beek and T. Fisher,
``Computing the Cassels--Tate pairing on 3-isogeny Selmer groups via cubic
norm equations,'' \emph{Acta Arithmetica} \textbf{185} (2018), 367--396,
\url{https://doi.org/10.4064/aa171108-11-4}.

\bibitem{MagmaHandbook}
W. Bosma, J. Cannon, and C. Playoust,
``The Magma algebra system. I. The user language,''
\emph{J. Symbolic Comput.} \textbf{24} (1997), 235--265;
current descent interface documented at
\url{https://magma.maths.usyd.edu.au/magma/handbook/text/1570}.
\end{thebibliography}
}

\end{document}
